Com esta pesquisa, tem-se como principal objetivo Desenvolver um sistema inteligente de monitoramento para o cultivo de cogumelos, utilizando sensores ambientais e inteligência artificial, a fim de prevenir infecções e infestações por parasitas, proporcionando maior controle, segurança e produtividade ao produtor, utilizando níveis de umidade, temperatura e índice de infecção. Além disso, como metas parciais visa-se:

- Realizar o treinamento do modelo de aprendizagem profunda (Deep Learning) visando alcançar uma acurácia mínima de 50\% e uma taxa de falsos positivos inferior a 10\%, assegurando desempenho satisfatório na detecção e classificação dos dados analisados. 

- Implementar uma integração eficiente entre o sistema proposto e a área de atuação dos produtores do Vale do Ribeira, ampliando o acesso e a disponibilidade de dados para o aprimoramento contínuo do modelo de inteligência artificial. 

- Desenvolver a integração de câmeras e sensores de temperatura e umidade em um ambiente de Internet das Coisas (IoT) funcional, possibilitando a coleta automatizada de dados e o monitoramento em tempo real das condições ambientais do cultivo. 